% =====================================================================
%  Ejercicio13PanquequesUnoDos.tex
%  ========================
%
%
%  Descripción:
%  ------------
%  Ejercicio 13 de la hoja de cotas asintóticas, recurrencias y
%  montículos, resuelto paso a paso. Un ejercicio por documento.
%
%
%  Considerations:
%  ------------
%  - Compila con pdflatex sin shell-escape. Dos pasadas, o latexmk
%  - Solo declara las convenciones que este ejercicio usa
%  - Sin título grande y sin pie de página: el encabezado de un
%    renglón deja la hoja entera para el desarrollo
%
%
%  Metadata:
%  ----------
%   * Author : zxxz6 (Bryan Violante Arriaga)
%   * Version: 1.0.0
%
%
%  History:
%  ------------
%  Author      Date            Description
%  zxxz6       24/09/2026      Condense la correctitud a dos paginas
%  zxxz6       24/09/2026      Creation
%
%
% =====================================================================
\documentclass[11pt,letterpaper]{article}

\newcommand{\materia}{Hoja de recurrencias y montículos}
\newcommand{\profesor}{Dr. José Martínez Carranza}
\newcommand{\periodo}{Otoño 2026}
\newcommand{\alumno}{Bryan Ezequiel Violante Arriaga}

% =====================================================================
%  Preambulo.tex
%  ========================
%
%
%  Descripción:
%  ------------
%  Preámbulo compartido de los apuntes de clase: paquetes, estilo de
%  página, entornos de teorema, las cuatro cajas de color, el estilo de
%  listings y el pseudocódigo en español. Se carga con \input desde
%  cada documento en lugar de copiarlo, para que un arreglo aquí llegue
%  a todas las libretas de una vez.
%
%
%  Considerations:
%  ------------
%  - Se carga con \input y ruta relativa, no con \usepackage, porque
%    \usepackage solo busca en la carpeta del documento y en el árbol
%    de TeX. Un .sty instalado en TEXMFHOME evitaría la ruta, pero vive
%    fuera del repo y una clona nueva no compilaría
%  - Los cuatro datos de la materia se declaran en el documento ANTES
%    del \input. hyperref expande \materia al cargarse, así que si no
%    existe todavía la compilación truena
%  - Cuando la materia la imparte más de un profesor, los nombres van
%    en un solo \profesor separados por \\. La portadilla los imprime
%    dentro de un center, que es donde \\ sí corta renglón. Repetir
%    \newcommand{\profesor} no es opción: truena con "already defined"
%  - algorithmicx no tiene opción de idioma y babel no lo traduce, así
%    que las palabras clave del pseudocódigo se renombran a mano
%  - Los nombres de los entornos van sin acento (\begin{definicion})
%    aunque el rótulo impreso sí lo lleve
%  - babel se carga con es-noshorthands porque en español el " queda
%    activo y se come el texto que le sigue, además de romper listings
%
%
%  Metadata:
%  ----------
%   * Author : zxxz6 (Bryan Violante Arriaga)
%   * Version: 1.7.0
%
%
%  History:
%  ------------
%  Author      Date            Description
%  zxxz6       24/09/2026      El nombre del alumno en el encabezado
%  zxxz6       24/09/2026      Encabezado de una hoja de un ejercicio
%  zxxz6       23/09/2026      Macros de los documentos de soluciones
%  zxxz6       30/08/2026      Macros de lim, limsup y liminf sobre n
%  zxxz6       25/08/2026      Macros de Omega, Theta, o y omega chicas
%  zxxz6       20/08/2026      Cargue pgfplots para graficas de datos
%  zxxz6       20/08/2026      Cargue tikz para diagramas de conjuntos
%  zxxz6       19/08/2026      Documente \profesor con varios nombres
%  zxxz6       19/08/2026      Agregue la caja Idea Random
%  zxxz6       18/08/2026      Creation
%
%
% =====================================================================
%  USO
%  ---
%  \documentclass[11pt,letterpaper]{article}
%  \newcommand{\materia}{...}      % los cuatro datos van arriba
%  \newcommand{\profesor}{...}
%  \newcommand{\periodo}{...}
%  \newcommand{\alumno}{...}
%  % =====================================================================
%  Preambulo.tex
%  ========================
%
%
%  Descripción:
%  ------------
%  Preámbulo compartido de los apuntes de clase: paquetes, estilo de
%  página, entornos de teorema, las cuatro cajas de color, el estilo de
%  listings y el pseudocódigo en español. Se carga con \input desde
%  cada documento en lugar de copiarlo, para que un arreglo aquí llegue
%  a todas las libretas de una vez.
%
%
%  Considerations:
%  ------------
%  - Se carga con \input y ruta relativa, no con \usepackage, porque
%    \usepackage solo busca en la carpeta del documento y en el árbol
%    de TeX. Un .sty instalado en TEXMFHOME evitaría la ruta, pero vive
%    fuera del repo y una clona nueva no compilaría
%  - Los cuatro datos de la materia se declaran en el documento ANTES
%    del \input. hyperref expande \materia al cargarse, así que si no
%    existe todavía la compilación truena
%  - Cuando la materia la imparte más de un profesor, los nombres van
%    en un solo \profesor separados por \\. La portadilla los imprime
%    dentro de un center, que es donde \\ sí corta renglón. Repetir
%    \newcommand{\profesor} no es opción: truena con "already defined"
%  - algorithmicx no tiene opción de idioma y babel no lo traduce, así
%    que las palabras clave del pseudocódigo se renombran a mano
%  - Los nombres de los entornos van sin acento (\begin{definicion})
%    aunque el rótulo impreso sí lo lleve
%  - babel se carga con es-noshorthands porque en español el " queda
%    activo y se come el texto que le sigue, además de romper listings
%
%
%  Metadata:
%  ----------
%   * Author : zxxz6 (Bryan Violante Arriaga)
%   * Version: 1.7.0
%
%
%  History:
%  ------------
%  Author      Date            Description
%  zxxz6       24/09/2026      El nombre del alumno en el encabezado
%  zxxz6       24/09/2026      Encabezado de una hoja de un ejercicio
%  zxxz6       23/09/2026      Macros de los documentos de soluciones
%  zxxz6       30/08/2026      Macros de lim, limsup y liminf sobre n
%  zxxz6       25/08/2026      Macros de Omega, Theta, o y omega chicas
%  zxxz6       20/08/2026      Cargue pgfplots para graficas de datos
%  zxxz6       20/08/2026      Cargue tikz para diagramas de conjuntos
%  zxxz6       19/08/2026      Documente \profesor con varios nombres
%  zxxz6       19/08/2026      Agregue la caja Idea Random
%  zxxz6       18/08/2026      Creation
%
%
% =====================================================================
%  USO
%  ---
%  \documentclass[11pt,letterpaper]{article}
%  \newcommand{\materia}{...}      % los cuatro datos van arriba
%  \newcommand{\profesor}{...}
%  \newcommand{\periodo}{...}
%  \newcommand{\alumno}{...}
%  % =====================================================================
%  Preambulo.tex
%  ========================
%
%
%  Descripción:
%  ------------
%  Preámbulo compartido de los apuntes de clase: paquetes, estilo de
%  página, entornos de teorema, las cuatro cajas de color, el estilo de
%  listings y el pseudocódigo en español. Se carga con \input desde
%  cada documento en lugar de copiarlo, para que un arreglo aquí llegue
%  a todas las libretas de una vez.
%
%
%  Considerations:
%  ------------
%  - Se carga con \input y ruta relativa, no con \usepackage, porque
%    \usepackage solo busca en la carpeta del documento y en el árbol
%    de TeX. Un .sty instalado en TEXMFHOME evitaría la ruta, pero vive
%    fuera del repo y una clona nueva no compilaría
%  - Los cuatro datos de la materia se declaran en el documento ANTES
%    del \input. hyperref expande \materia al cargarse, así que si no
%    existe todavía la compilación truena
%  - Cuando la materia la imparte más de un profesor, los nombres van
%    en un solo \profesor separados por \\. La portadilla los imprime
%    dentro de un center, que es donde \\ sí corta renglón. Repetir
%    \newcommand{\profesor} no es opción: truena con "already defined"
%  - algorithmicx no tiene opción de idioma y babel no lo traduce, así
%    que las palabras clave del pseudocódigo se renombran a mano
%  - Los nombres de los entornos van sin acento (\begin{definicion})
%    aunque el rótulo impreso sí lo lleve
%  - babel se carga con es-noshorthands porque en español el " queda
%    activo y se come el texto que le sigue, además de romper listings
%
%
%  Metadata:
%  ----------
%   * Author : zxxz6 (Bryan Violante Arriaga)
%   * Version: 1.7.0
%
%
%  History:
%  ------------
%  Author      Date            Description
%  zxxz6       24/09/2026      El nombre del alumno en el encabezado
%  zxxz6       24/09/2026      Encabezado de una hoja de un ejercicio
%  zxxz6       23/09/2026      Macros de los documentos de soluciones
%  zxxz6       30/08/2026      Macros de lim, limsup y liminf sobre n
%  zxxz6       25/08/2026      Macros de Omega, Theta, o y omega chicas
%  zxxz6       20/08/2026      Cargue pgfplots para graficas de datos
%  zxxz6       20/08/2026      Cargue tikz para diagramas de conjuntos
%  zxxz6       19/08/2026      Documente \profesor con varios nombres
%  zxxz6       19/08/2026      Agregue la caja Idea Random
%  zxxz6       18/08/2026      Creation
%
%
% =====================================================================
%  USO
%  ---
%  \documentclass[11pt,letterpaper]{article}
%  \newcommand{\materia}{...}      % los cuatro datos van arriba
%  \newcommand{\profesor}{...}
%  \newcommand{\periodo}{...}
%  \newcommand{\alumno}{...}
%  \input{../../templates/Preambulo.tex}
%  \begin{document}
%  \portadilla
%
%  Con dos o mas profesores se declara UN solo \profesor y los nombres
%  se separan con \\, uno por renglon:
%
%  \newcommand{\profesor}{Dra. Fulana de Tal \\
%                         Dr. Mengano de Cual}
% =====================================================================


% ==================== IDIOMA Y CODIFICACIÓN ====================
\usepackage[utf8]{inputenc}
\usepackage[T1]{fontenc}
\usepackage[spanish,es-tabla,es-noshorthands]{babel}
\usepackage{lmodern}
\usepackage{microtype}

% ==================== PÁGINA ====================
\usepackage[letterpaper,margin=2.3cm,headheight=14pt]{geometry}
\usepackage{parskip}
\usepackage{fancyhdr}
\usepackage{enumitem}

% ==================== MATEMÁTICAS ====================
\usepackage{amsmath,amssymb,amsthm,mathtools}

% ==================== FIGURAS, TABLAS, CÓDIGO ====================
\usepackage{graphicx}
\usepackage{booktabs}
\usepackage{array}
% float da el especificador [H], que clava una tabla donde se
% escribio. Una traza que LaTeX manda dos paginas adelante deja
% de leerse junto al paso que la produjo
\usepackage{float}
\usepackage{xcolor}
% tikz ya viene arrastrado por tcolorbox[most], pero se carga
% explicito para no depender de ese detalle; la libreria babel
% protege a tikz de los caracteres activos del español
\usepackage{tikz}
\usetikzlibrary{babel}
% pgfplots dibuja graficas de datos (ejes, escalas log) sobre tikz
\usepackage{pgfplots}
\pgfplotsset{compat=1.18}
\usepackage{listings}
\usepackage[ruled,section]{algorithm}
\usepackage{algpseudocode}
\usepackage[most]{tcolorbox}
% soul da \hl, el unico resaltado que parte renglon; \colorbox no
\usepackage{soul}
\usepackage{hyperref}

% =====================================================================
%  ESTILO
% =====================================================================

% --- encabezado y pie de cada página ---
\pagestyle{fancy}
\fancyhf{}
\fancyhead[L]{\small\textsc{\materia}}
\fancyhead[R]{\small\periodo}
\fancyfoot[C]{\small\thepage}
\renewcommand{\headrulewidth}{0.4pt}

% --- enlaces ---
\hypersetup{
  colorlinks=true,
  linkcolor=black,
  urlcolor=blue,
  citecolor=black,
  pdftitle={Apuntes de \materia},
  pdfauthor={\alumno}
}

% --- paleta ---
\definecolor{acento}{HTML}{1F4E79}
\definecolor{fondoIdea}{HTML}{EAF2FB}
\definecolor{fondoOjo}{HTML}{FDF0E6}
\definecolor{fondoPend}{HTML}{F2F0F7}
\definecolor{comentario}{HTML}{5A7D5A}
\definecolor{cadena}{HTML}{A03030}
\definecolor{fondoCodigo}{HTML}{F7F7F7}
\definecolor{fondoResultado}{HTML}{EAF5EC}
\definecolor{marcador}{HTML}{FFF3A3}

% --- entornos de teorema; numeran por sección, o sea por clase ---
\theoremstyle{definition}
\newtheorem{definicion}{Definición}[section]
\newtheorem{ejemplo}[definicion]{Ejemplo}
\newtheorem{ejercicio}[definicion]{Ejercicio}

\theoremstyle{plain}
\newtheorem{teorema}[definicion]{Teorema}
\newtheorem{lema}[definicion]{Lema}
\newtheorem{corolario}[definicion]{Corolario}
\newtheorem{proposicion}[definicion]{Proposición}

\theoremstyle{remark}
\newtheorem*{observacion}{Observación}
\newtheorem*{quick_sum}{Resumen rápido}

% --- cajas de color para lo que se relee antes del examen ---
\newtcolorbox{idea}[1][]{
  colback=fondoIdea, colframe=acento, boxrule=0.6pt,
  arc=2pt, left=6pt, right=6pt, top=5pt, bottom=5pt,
  title={Idea clave}, fonttitle=\bfseries\small, #1
}
\newtcolorbox{ojo}[1][]{
  colback=fondoOjo, colframe=orange!70!black, boxrule=0.6pt,
  arc=2pt, left=6pt, right=6pt, top=5pt, bottom=5pt,
  title={Ojito...}, fonttitle=\bfseries\small, #1
}
\newtcolorbox{pendiente}[1][]{
  colback=fondoPend, colframe=violet!60!black, boxrule=0.6pt,
  arc=2pt, left=6pt, right=6pt, top=5pt, bottom=5pt,
  title={Pendiente por resolver}, fonttitle=\bfseries\small, #1
}
\newtcolorbox{idear}[1][]{
  colback=fondoPend, colframe=pink!60!black, boxrule=0.6pt,
  arc=2pt, left=6pt, right=6pt, top=5pt, bottom=5pt,
  title={Idea Random}, fonttitle=\bfseries\small, #1
}

% --- el resultado al que llego un ejercicio, para que se vea de
%     lejos al hojear la hoja antes del examen ---
\newtcolorbox{resultado}[1][]{
  colback=fondoResultado, colframe=green!45!black, boxrule=0.6pt,
  arc=2pt, left=6pt, right=6pt, top=5pt, bottom=5pt,
  title={Resultado}, fonttitle=\bfseries\small, #1
}

% --- código ---
\lstdefinestyle{codigo}{
  backgroundcolor=\color{fondoCodigo},
  basicstyle=\ttfamily\footnotesize,
  keywordstyle=\color{acento}\bfseries,
  commentstyle=\color{comentario}\itshape,
  stringstyle=\color{cadena},
  numbers=left, numberstyle=\tiny\color{gray}, numbersep=8pt,
  frame=single, rulecolor=\color{gray!40},
  breaklines=true, showstringspaces=false,
  tabsize=4, captionpos=b,
  extendedchars=true, inputencoding=utf8,
  literate={á}{{\'a}}1 {é}{{\'e}}1 {í}{{\'i}}1
           {ó}{{\'o}}1 {ú}{{\'u}}1 {ñ}{{\~n}}1
           {Ü}{{\"U}}1 {ü}{{\"u}}1 {¿}{{?`}}1
           {¡}{{!`}}1
}
\lstset{style=codigo}
\renewcommand{\lstlistingname}{Código}
\renewcommand{\lstlistlistingname}{Índice de código}
% --- pseudocódigo; algorithmicx arma el cuerpo y algorithm el
%     flotante que lo numera y lo pone en el índice. Las palabras
%     clave vienen en inglés de fábrica y babel no las toca, así que
%     se renombran una por una ---
\floatname{algorithm}{Algoritmo}
\renewcommand{\listalgorithmname}{Índice de algoritmos}

\algrenewcommand\algorithmicrequire{\textbf{Entrada:}}
\algrenewcommand\algorithmicensure{\textbf{Salida:}}
\algrenewcommand\algorithmicend{\textbf{fin}}
\algrenewcommand\algorithmicif{\textbf{si}}
\algrenewcommand\algorithmicthen{\textbf{entonces}}
\algrenewcommand\algorithmicelse{\textbf{si no}}
\algrenewcommand\algorithmicfor{\textbf{para}}
\algrenewcommand\algorithmicforall{\textbf{para todo}}
\algrenewcommand\algorithmicdo{\textbf{hacer}}
\algrenewcommand\algorithmicwhile{\textbf{mientras}}
\algrenewcommand\algorithmicrepeat{\textbf{repetir}}
\algrenewcommand\algorithmicuntil{\textbf{hasta que}}
\algrenewcommand\algorithmicloop{\textbf{ciclo}}
\algrenewcommand\algorithmicprocedure{\textbf{procedimiento}}
\algrenewcommand\algorithmicfunction{\textbf{función}}
\algrenewcommand\algorithmicreturn{\textbf{devolver}}
\algrenewcomment[1]{\hfill{\color{comentario}\itshape // #1}}

% Las dos que algorithmicx no trae. El \ del final es obligatorio: sin
% el TeX se come el espacio y sale "hastan-1" en vez de "hasta n-1"
\algnewcommand\Hasta{\textbf{hasta}\ }
\algnewcommand\Intercambiar{\textbf{intercambiar}\ }

% --- abreviaturas que se usan a cada rato ---
\newcommand{\N}{\mathbb{N}}
\newcommand{\Z}{\mathbb{Z}}
\newcommand{\Q}{\mathbb{Q}}
\newcommand{\R}{\mathbb{R}}
\newcommand{\C}{\mathbb{C}}
\newcommand{\abs}[1]{\left\lvert #1 \right\rvert}
\newcommand{\norm}[1]{\left\lVert #1 \right\rVert}
\newcommand{\set}[1]{\left\{ #1 \right\}}
\newcommand{\Ord}[1]{\mathcal{O}\!\left(#1\right)}
% Mayuscula la notacion grande, minuscula la chica
\newcommand{\Omg}[1]{\Omega\!\left(#1\right)}
\newcommand{\Tht}[1]{\Theta\!\left(#1\right)}
\newcommand{\ord}[1]{o\!\left(#1\right)}
\newcommand{\omg}[1]{\omega\!\left(#1\right)}
% Los tres limites sobre n, que salen en cada renglon del criterio
% del cociente para la notacion asintotica
\newcommand{\LM}{\lim_{n \to \infty}}
\newcommand{\LS}{\limsup_{n \to \infty}}
\newcommand{\LI}{\liminf_{n \to \infty}}

% --- encabezado de cada sesión: título a la izquierda, fecha a la
%     derecha. El argumento corto es el que sale en el índice ---
\newcommand{\clase}[2]{%
  \section[#1]{#1 \hfill \normalfont\small\textit{#2}}%
}

% --- portadilla y tabla de contenidos. Se llama una vez, justo
%     despues de \begin{document}. El renglon del profesor va dentro
%     del center, asi que \profesor puede traer varios nombres
%     separados por \\ y salen apilados y centrados sin tocar nada ---
\newcommand{\portadilla}{%
  \begin{center}
    {\LARGE\bfseries\materia\par}
    \vspace{0.4em}
    {\large \profesor \par}
    \vspace{0.2em}
    {\small \periodo\quad-\quad\alumno\par}
    \vspace{0.3em}
    \rule{\linewidth}{0.5pt}
  \end{center}
  \vspace{0.5em}
  \tableofcontents
  \vspace{1em}
  \rule{\linewidth}{0.5pt}%
}

% =====================================================================
%  DOCUMENTOS DE SOLUCIONES
%  Lo que usa una tanda de ejercicios resueltos. Una libreta no toca
%  nada de aqui, pero vive en el preambulo compartido para no tener
%  dos archivos que arreglar cuando cambie el estilo.
% =====================================================================

% --- encabezado compacto: titulo, subtitulo y regla. Sustituye a
%     \portadilla cuando el documento es una tanda de ejercicios: sin
%     indice, porque se hojea de corrido, no se navega ---
\newcommand{\encabezadoSoluciones}[2]{%
  % El paquete algorithm se carga con la opcion [section] para que en
  % una libreta el pseudocodigo se numere por clase. Aqui no hay
  % secciones, asi que esa numeracion sale como "Algoritmo 0.1". Se
  % corrige al vuelo, y solo en los documentos que llaman a este
  % encabezado: la libreta no se entera.
  \renewcommand{\thealgorithm}{\arabic{algorithm}}%
  \begin{center}
    {\LARGE\bfseries #1\par}
    \vspace{0.3em}
    {\small\itshape #2\par}
    \vspace{0.5em}
    \rule{\linewidth}{0.5pt}
  \end{center}
  \vspace{0.3em}%
}

% --- hoja de un solo ejercicio: un renglon chico arriba con el
%     titulo y su numero, y ya. Sin titulo grande y sin pie de pagina,
%     porque en una hoja de dos paginas cada centimetro que se va en
%     adornos es un paso del desarrollo que no cabe ---
\newcommand{\encabezadoEjercicio}[3]{%
  % Misma correccion que en \encabezadoSoluciones: sin secciones, el
  % pseudocodigo saldria numerado "0.1"
  \renewcommand{\thealgorithm}{\arabic{algorithm}}%
  \fancyhf{}%
  \fancyhead[L]{\small\textsc{#1} (ejercicio #2 de #3)}%
  % El nombre va a la derecha porque las hojas de respuestas se
  % entregan y el profesor pide que cada una lo lleve
  \fancyhead[R]{\small\alumno}%
  \renewcommand{\headrulewidth}{0.4pt}%
  \pagestyle{fancy}%
  \thispagestyle{fancy}%
}

% --- bloque tematico, numerado con romanos: I, II, III. Agrupa los
%     ejercicios que comparten tema. No dibuja regla: la del
%     \problema que sigue ya separa el titulo del primer ejercicio,
%     y dos rayas seguidas dejan un hueco feo ---
\newcounter{parte}
\renewcommand{\theparte}{\Roman{parte}}
\newcommand{\parte}[1]{%
  \bigskip
  \refstepcounter{parte}%
  \noindent{\large\bfseries \theparte. #1\par}%
}

% --- un ejercicio resuelto. Se numera solo y corrido a lo largo del
%     documento, sin reiniciar en cada parte, para poder decir "el 14"
%     sin aclarar de que bloque ---
\newcounter{problema}
\newcommand{\problema}[1]{%
  \bigskip
  \noindent\rule{\linewidth}{0.4pt}\par
  \vspace{0.3em}
  \refstepcounter{problema}%
  \noindent{\bfseries Ejercicio \theproblema\ --- #1\par}
  \vspace{0.3em}%
}

% --- rotulo del paso que sigue: Caso base, Hipotesis inductiva, Paso
%     inductivo. El texto sigue en el mismo renglon ---
\newcommand{\paso}[1]{\par\vspace{0.3em}\noindent\textbf{#1}\ }

% --- la respuesta de un ejercicio, centrada y en un cuadro. El
%     argumento va en modo matematico, sin los delimitadores ---
\newcommand{\respuesta}[1]{%
  \[
    \boxed{\;#1\;}
  \]
}

% --- resaltado inline, como pasarle el plumon a media frase. No
%     admite matematicas adentro: soul rompe con $...$, asi que una
%     formula que hay que destacar va en \respuesta ---
\sethlcolor{marcador}
\newcommand{\resaltar}[1]{\hl{#1}}

% --- justificacion al final de un renglon de \align, en chiquito y
%     entre parentesis, para decir de donde salio el paso ---
\newcommand{\razon}[1]{\quad\text{\small\itshape (#1)}}

%############################ END OF PREAMBULO.TEX #############################
%###############################################################################

%  \begin{document}
%  \portadilla
%
%  Con dos o mas profesores se declara UN solo \profesor y los nombres
%  se separan con \\, uno por renglon:
%
%  \newcommand{\profesor}{Dra. Fulana de Tal \\
%                         Dr. Mengano de Cual}
% =====================================================================


% ==================== IDIOMA Y CODIFICACIÓN ====================
\usepackage[utf8]{inputenc}
\usepackage[T1]{fontenc}
\usepackage[spanish,es-tabla,es-noshorthands]{babel}
\usepackage{lmodern}
\usepackage{microtype}

% ==================== PÁGINA ====================
\usepackage[letterpaper,margin=2.3cm,headheight=14pt]{geometry}
\usepackage{parskip}
\usepackage{fancyhdr}
\usepackage{enumitem}

% ==================== MATEMÁTICAS ====================
\usepackage{amsmath,amssymb,amsthm,mathtools}

% ==================== FIGURAS, TABLAS, CÓDIGO ====================
\usepackage{graphicx}
\usepackage{booktabs}
\usepackage{array}
% float da el especificador [H], que clava una tabla donde se
% escribio. Una traza que LaTeX manda dos paginas adelante deja
% de leerse junto al paso que la produjo
\usepackage{float}
\usepackage{xcolor}
% tikz ya viene arrastrado por tcolorbox[most], pero se carga
% explicito para no depender de ese detalle; la libreria babel
% protege a tikz de los caracteres activos del español
\usepackage{tikz}
\usetikzlibrary{babel}
% pgfplots dibuja graficas de datos (ejes, escalas log) sobre tikz
\usepackage{pgfplots}
\pgfplotsset{compat=1.18}
\usepackage{listings}
\usepackage[ruled,section]{algorithm}
\usepackage{algpseudocode}
\usepackage[most]{tcolorbox}
% soul da \hl, el unico resaltado que parte renglon; \colorbox no
\usepackage{soul}
\usepackage{hyperref}

% =====================================================================
%  ESTILO
% =====================================================================

% --- encabezado y pie de cada página ---
\pagestyle{fancy}
\fancyhf{}
\fancyhead[L]{\small\textsc{\materia}}
\fancyhead[R]{\small\periodo}
\fancyfoot[C]{\small\thepage}
\renewcommand{\headrulewidth}{0.4pt}

% --- enlaces ---
\hypersetup{
  colorlinks=true,
  linkcolor=black,
  urlcolor=blue,
  citecolor=black,
  pdftitle={Apuntes de \materia},
  pdfauthor={\alumno}
}

% --- paleta ---
\definecolor{acento}{HTML}{1F4E79}
\definecolor{fondoIdea}{HTML}{EAF2FB}
\definecolor{fondoOjo}{HTML}{FDF0E6}
\definecolor{fondoPend}{HTML}{F2F0F7}
\definecolor{comentario}{HTML}{5A7D5A}
\definecolor{cadena}{HTML}{A03030}
\definecolor{fondoCodigo}{HTML}{F7F7F7}
\definecolor{fondoResultado}{HTML}{EAF5EC}
\definecolor{marcador}{HTML}{FFF3A3}

% --- entornos de teorema; numeran por sección, o sea por clase ---
\theoremstyle{definition}
\newtheorem{definicion}{Definición}[section]
\newtheorem{ejemplo}[definicion]{Ejemplo}
\newtheorem{ejercicio}[definicion]{Ejercicio}

\theoremstyle{plain}
\newtheorem{teorema}[definicion]{Teorema}
\newtheorem{lema}[definicion]{Lema}
\newtheorem{corolario}[definicion]{Corolario}
\newtheorem{proposicion}[definicion]{Proposición}

\theoremstyle{remark}
\newtheorem*{observacion}{Observación}
\newtheorem*{quick_sum}{Resumen rápido}

% --- cajas de color para lo que se relee antes del examen ---
\newtcolorbox{idea}[1][]{
  colback=fondoIdea, colframe=acento, boxrule=0.6pt,
  arc=2pt, left=6pt, right=6pt, top=5pt, bottom=5pt,
  title={Idea clave}, fonttitle=\bfseries\small, #1
}
\newtcolorbox{ojo}[1][]{
  colback=fondoOjo, colframe=orange!70!black, boxrule=0.6pt,
  arc=2pt, left=6pt, right=6pt, top=5pt, bottom=5pt,
  title={Ojito...}, fonttitle=\bfseries\small, #1
}
\newtcolorbox{pendiente}[1][]{
  colback=fondoPend, colframe=violet!60!black, boxrule=0.6pt,
  arc=2pt, left=6pt, right=6pt, top=5pt, bottom=5pt,
  title={Pendiente por resolver}, fonttitle=\bfseries\small, #1
}
\newtcolorbox{idear}[1][]{
  colback=fondoPend, colframe=pink!60!black, boxrule=0.6pt,
  arc=2pt, left=6pt, right=6pt, top=5pt, bottom=5pt,
  title={Idea Random}, fonttitle=\bfseries\small, #1
}

% --- el resultado al que llego un ejercicio, para que se vea de
%     lejos al hojear la hoja antes del examen ---
\newtcolorbox{resultado}[1][]{
  colback=fondoResultado, colframe=green!45!black, boxrule=0.6pt,
  arc=2pt, left=6pt, right=6pt, top=5pt, bottom=5pt,
  title={Resultado}, fonttitle=\bfseries\small, #1
}

% --- código ---
\lstdefinestyle{codigo}{
  backgroundcolor=\color{fondoCodigo},
  basicstyle=\ttfamily\footnotesize,
  keywordstyle=\color{acento}\bfseries,
  commentstyle=\color{comentario}\itshape,
  stringstyle=\color{cadena},
  numbers=left, numberstyle=\tiny\color{gray}, numbersep=8pt,
  frame=single, rulecolor=\color{gray!40},
  breaklines=true, showstringspaces=false,
  tabsize=4, captionpos=b,
  extendedchars=true, inputencoding=utf8,
  literate={á}{{\'a}}1 {é}{{\'e}}1 {í}{{\'i}}1
           {ó}{{\'o}}1 {ú}{{\'u}}1 {ñ}{{\~n}}1
           {Ü}{{\"U}}1 {ü}{{\"u}}1 {¿}{{?`}}1
           {¡}{{!`}}1
}
\lstset{style=codigo}
\renewcommand{\lstlistingname}{Código}
\renewcommand{\lstlistlistingname}{Índice de código}
% --- pseudocódigo; algorithmicx arma el cuerpo y algorithm el
%     flotante que lo numera y lo pone en el índice. Las palabras
%     clave vienen en inglés de fábrica y babel no las toca, así que
%     se renombran una por una ---
\floatname{algorithm}{Algoritmo}
\renewcommand{\listalgorithmname}{Índice de algoritmos}

\algrenewcommand\algorithmicrequire{\textbf{Entrada:}}
\algrenewcommand\algorithmicensure{\textbf{Salida:}}
\algrenewcommand\algorithmicend{\textbf{fin}}
\algrenewcommand\algorithmicif{\textbf{si}}
\algrenewcommand\algorithmicthen{\textbf{entonces}}
\algrenewcommand\algorithmicelse{\textbf{si no}}
\algrenewcommand\algorithmicfor{\textbf{para}}
\algrenewcommand\algorithmicforall{\textbf{para todo}}
\algrenewcommand\algorithmicdo{\textbf{hacer}}
\algrenewcommand\algorithmicwhile{\textbf{mientras}}
\algrenewcommand\algorithmicrepeat{\textbf{repetir}}
\algrenewcommand\algorithmicuntil{\textbf{hasta que}}
\algrenewcommand\algorithmicloop{\textbf{ciclo}}
\algrenewcommand\algorithmicprocedure{\textbf{procedimiento}}
\algrenewcommand\algorithmicfunction{\textbf{función}}
\algrenewcommand\algorithmicreturn{\textbf{devolver}}
\algrenewcomment[1]{\hfill{\color{comentario}\itshape // #1}}

% Las dos que algorithmicx no trae. El \ del final es obligatorio: sin
% el TeX se come el espacio y sale "hastan-1" en vez de "hasta n-1"
\algnewcommand\Hasta{\textbf{hasta}\ }
\algnewcommand\Intercambiar{\textbf{intercambiar}\ }

% --- abreviaturas que se usan a cada rato ---
\newcommand{\N}{\mathbb{N}}
\newcommand{\Z}{\mathbb{Z}}
\newcommand{\Q}{\mathbb{Q}}
\newcommand{\R}{\mathbb{R}}
\newcommand{\C}{\mathbb{C}}
\newcommand{\abs}[1]{\left\lvert #1 \right\rvert}
\newcommand{\norm}[1]{\left\lVert #1 \right\rVert}
\newcommand{\set}[1]{\left\{ #1 \right\}}
\newcommand{\Ord}[1]{\mathcal{O}\!\left(#1\right)}
% Mayuscula la notacion grande, minuscula la chica
\newcommand{\Omg}[1]{\Omega\!\left(#1\right)}
\newcommand{\Tht}[1]{\Theta\!\left(#1\right)}
\newcommand{\ord}[1]{o\!\left(#1\right)}
\newcommand{\omg}[1]{\omega\!\left(#1\right)}
% Los tres limites sobre n, que salen en cada renglon del criterio
% del cociente para la notacion asintotica
\newcommand{\LM}{\lim_{n \to \infty}}
\newcommand{\LS}{\limsup_{n \to \infty}}
\newcommand{\LI}{\liminf_{n \to \infty}}

% --- encabezado de cada sesión: título a la izquierda, fecha a la
%     derecha. El argumento corto es el que sale en el índice ---
\newcommand{\clase}[2]{%
  \section[#1]{#1 \hfill \normalfont\small\textit{#2}}%
}

% --- portadilla y tabla de contenidos. Se llama una vez, justo
%     despues de \begin{document}. El renglon del profesor va dentro
%     del center, asi que \profesor puede traer varios nombres
%     separados por \\ y salen apilados y centrados sin tocar nada ---
\newcommand{\portadilla}{%
  \begin{center}
    {\LARGE\bfseries\materia\par}
    \vspace{0.4em}
    {\large \profesor \par}
    \vspace{0.2em}
    {\small \periodo\quad-\quad\alumno\par}
    \vspace{0.3em}
    \rule{\linewidth}{0.5pt}
  \end{center}
  \vspace{0.5em}
  \tableofcontents
  \vspace{1em}
  \rule{\linewidth}{0.5pt}%
}

% =====================================================================
%  DOCUMENTOS DE SOLUCIONES
%  Lo que usa una tanda de ejercicios resueltos. Una libreta no toca
%  nada de aqui, pero vive en el preambulo compartido para no tener
%  dos archivos que arreglar cuando cambie el estilo.
% =====================================================================

% --- encabezado compacto: titulo, subtitulo y regla. Sustituye a
%     \portadilla cuando el documento es una tanda de ejercicios: sin
%     indice, porque se hojea de corrido, no se navega ---
\newcommand{\encabezadoSoluciones}[2]{%
  % El paquete algorithm se carga con la opcion [section] para que en
  % una libreta el pseudocodigo se numere por clase. Aqui no hay
  % secciones, asi que esa numeracion sale como "Algoritmo 0.1". Se
  % corrige al vuelo, y solo en los documentos que llaman a este
  % encabezado: la libreta no se entera.
  \renewcommand{\thealgorithm}{\arabic{algorithm}}%
  \begin{center}
    {\LARGE\bfseries #1\par}
    \vspace{0.3em}
    {\small\itshape #2\par}
    \vspace{0.5em}
    \rule{\linewidth}{0.5pt}
  \end{center}
  \vspace{0.3em}%
}

% --- hoja de un solo ejercicio: un renglon chico arriba con el
%     titulo y su numero, y ya. Sin titulo grande y sin pie de pagina,
%     porque en una hoja de dos paginas cada centimetro que se va en
%     adornos es un paso del desarrollo que no cabe ---
\newcommand{\encabezadoEjercicio}[3]{%
  % Misma correccion que en \encabezadoSoluciones: sin secciones, el
  % pseudocodigo saldria numerado "0.1"
  \renewcommand{\thealgorithm}{\arabic{algorithm}}%
  \fancyhf{}%
  \fancyhead[L]{\small\textsc{#1} (ejercicio #2 de #3)}%
  % El nombre va a la derecha porque las hojas de respuestas se
  % entregan y el profesor pide que cada una lo lleve
  \fancyhead[R]{\small\alumno}%
  \renewcommand{\headrulewidth}{0.4pt}%
  \pagestyle{fancy}%
  \thispagestyle{fancy}%
}

% --- bloque tematico, numerado con romanos: I, II, III. Agrupa los
%     ejercicios que comparten tema. No dibuja regla: la del
%     \problema que sigue ya separa el titulo del primer ejercicio,
%     y dos rayas seguidas dejan un hueco feo ---
\newcounter{parte}
\renewcommand{\theparte}{\Roman{parte}}
\newcommand{\parte}[1]{%
  \bigskip
  \refstepcounter{parte}%
  \noindent{\large\bfseries \theparte. #1\par}%
}

% --- un ejercicio resuelto. Se numera solo y corrido a lo largo del
%     documento, sin reiniciar en cada parte, para poder decir "el 14"
%     sin aclarar de que bloque ---
\newcounter{problema}
\newcommand{\problema}[1]{%
  \bigskip
  \noindent\rule{\linewidth}{0.4pt}\par
  \vspace{0.3em}
  \refstepcounter{problema}%
  \noindent{\bfseries Ejercicio \theproblema\ --- #1\par}
  \vspace{0.3em}%
}

% --- rotulo del paso que sigue: Caso base, Hipotesis inductiva, Paso
%     inductivo. El texto sigue en el mismo renglon ---
\newcommand{\paso}[1]{\par\vspace{0.3em}\noindent\textbf{#1}\ }

% --- la respuesta de un ejercicio, centrada y en un cuadro. El
%     argumento va en modo matematico, sin los delimitadores ---
\newcommand{\respuesta}[1]{%
  \[
    \boxed{\;#1\;}
  \]
}

% --- resaltado inline, como pasarle el plumon a media frase. No
%     admite matematicas adentro: soul rompe con $...$, asi que una
%     formula que hay que destacar va en \respuesta ---
\sethlcolor{marcador}
\newcommand{\resaltar}[1]{\hl{#1}}

% --- justificacion al final de un renglon de \align, en chiquito y
%     entre parentesis, para decir de donde salio el paso ---
\newcommand{\razon}[1]{\quad\text{\small\itshape (#1)}}

%############################ END OF PREAMBULO.TEX #############################
%###############################################################################

%  \begin{document}
%  \portadilla
%
%  Con dos o mas profesores se declara UN solo \profesor y los nombres
%  se separan con \\, uno por renglon:
%
%  \newcommand{\profesor}{Dra. Fulana de Tal \\
%                         Dr. Mengano de Cual}
% =====================================================================


% ==================== IDIOMA Y CODIFICACIÓN ====================
\usepackage[utf8]{inputenc}
\usepackage[T1]{fontenc}
\usepackage[spanish,es-tabla,es-noshorthands]{babel}
\usepackage{lmodern}
\usepackage{microtype}

% ==================== PÁGINA ====================
\usepackage[letterpaper,margin=2.3cm,headheight=14pt]{geometry}
\usepackage{parskip}
\usepackage{fancyhdr}
\usepackage{enumitem}

% ==================== MATEMÁTICAS ====================
\usepackage{amsmath,amssymb,amsthm,mathtools}

% ==================== FIGURAS, TABLAS, CÓDIGO ====================
\usepackage{graphicx}
\usepackage{booktabs}
\usepackage{array}
% float da el especificador [H], que clava una tabla donde se
% escribio. Una traza que LaTeX manda dos paginas adelante deja
% de leerse junto al paso que la produjo
\usepackage{float}
\usepackage{xcolor}
% tikz ya viene arrastrado por tcolorbox[most], pero se carga
% explicito para no depender de ese detalle; la libreria babel
% protege a tikz de los caracteres activos del español
\usepackage{tikz}
\usetikzlibrary{babel}
% pgfplots dibuja graficas de datos (ejes, escalas log) sobre tikz
\usepackage{pgfplots}
\pgfplotsset{compat=1.18}
\usepackage{listings}
\usepackage[ruled,section]{algorithm}
\usepackage{algpseudocode}
\usepackage[most]{tcolorbox}
% soul da \hl, el unico resaltado que parte renglon; \colorbox no
\usepackage{soul}
\usepackage{hyperref}

% =====================================================================
%  ESTILO
% =====================================================================

% --- encabezado y pie de cada página ---
\pagestyle{fancy}
\fancyhf{}
\fancyhead[L]{\small\textsc{\materia}}
\fancyhead[R]{\small\periodo}
\fancyfoot[C]{\small\thepage}
\renewcommand{\headrulewidth}{0.4pt}

% --- enlaces ---
\hypersetup{
  colorlinks=true,
  linkcolor=black,
  urlcolor=blue,
  citecolor=black,
  pdftitle={Apuntes de \materia},
  pdfauthor={\alumno}
}

% --- paleta ---
\definecolor{acento}{HTML}{1F4E79}
\definecolor{fondoIdea}{HTML}{EAF2FB}
\definecolor{fondoOjo}{HTML}{FDF0E6}
\definecolor{fondoPend}{HTML}{F2F0F7}
\definecolor{comentario}{HTML}{5A7D5A}
\definecolor{cadena}{HTML}{A03030}
\definecolor{fondoCodigo}{HTML}{F7F7F7}
\definecolor{fondoResultado}{HTML}{EAF5EC}
\definecolor{marcador}{HTML}{FFF3A3}

% --- entornos de teorema; numeran por sección, o sea por clase ---
\theoremstyle{definition}
\newtheorem{definicion}{Definición}[section]
\newtheorem{ejemplo}[definicion]{Ejemplo}
\newtheorem{ejercicio}[definicion]{Ejercicio}

\theoremstyle{plain}
\newtheorem{teorema}[definicion]{Teorema}
\newtheorem{lema}[definicion]{Lema}
\newtheorem{corolario}[definicion]{Corolario}
\newtheorem{proposicion}[definicion]{Proposición}

\theoremstyle{remark}
\newtheorem*{observacion}{Observación}
\newtheorem*{quick_sum}{Resumen rápido}

% --- cajas de color para lo que se relee antes del examen ---
\newtcolorbox{idea}[1][]{
  colback=fondoIdea, colframe=acento, boxrule=0.6pt,
  arc=2pt, left=6pt, right=6pt, top=5pt, bottom=5pt,
  title={Idea clave}, fonttitle=\bfseries\small, #1
}
\newtcolorbox{ojo}[1][]{
  colback=fondoOjo, colframe=orange!70!black, boxrule=0.6pt,
  arc=2pt, left=6pt, right=6pt, top=5pt, bottom=5pt,
  title={Ojito...}, fonttitle=\bfseries\small, #1
}
\newtcolorbox{pendiente}[1][]{
  colback=fondoPend, colframe=violet!60!black, boxrule=0.6pt,
  arc=2pt, left=6pt, right=6pt, top=5pt, bottom=5pt,
  title={Pendiente por resolver}, fonttitle=\bfseries\small, #1
}
\newtcolorbox{idear}[1][]{
  colback=fondoPend, colframe=pink!60!black, boxrule=0.6pt,
  arc=2pt, left=6pt, right=6pt, top=5pt, bottom=5pt,
  title={Idea Random}, fonttitle=\bfseries\small, #1
}

% --- el resultado al que llego un ejercicio, para que se vea de
%     lejos al hojear la hoja antes del examen ---
\newtcolorbox{resultado}[1][]{
  colback=fondoResultado, colframe=green!45!black, boxrule=0.6pt,
  arc=2pt, left=6pt, right=6pt, top=5pt, bottom=5pt,
  title={Resultado}, fonttitle=\bfseries\small, #1
}

% --- código ---
\lstdefinestyle{codigo}{
  backgroundcolor=\color{fondoCodigo},
  basicstyle=\ttfamily\footnotesize,
  keywordstyle=\color{acento}\bfseries,
  commentstyle=\color{comentario}\itshape,
  stringstyle=\color{cadena},
  numbers=left, numberstyle=\tiny\color{gray}, numbersep=8pt,
  frame=single, rulecolor=\color{gray!40},
  breaklines=true, showstringspaces=false,
  tabsize=4, captionpos=b,
  extendedchars=true, inputencoding=utf8,
  literate={á}{{\'a}}1 {é}{{\'e}}1 {í}{{\'i}}1
           {ó}{{\'o}}1 {ú}{{\'u}}1 {ñ}{{\~n}}1
           {Ü}{{\"U}}1 {ü}{{\"u}}1 {¿}{{?`}}1
           {¡}{{!`}}1
}
\lstset{style=codigo}
\renewcommand{\lstlistingname}{Código}
\renewcommand{\lstlistlistingname}{Índice de código}
% --- pseudocódigo; algorithmicx arma el cuerpo y algorithm el
%     flotante que lo numera y lo pone en el índice. Las palabras
%     clave vienen en inglés de fábrica y babel no las toca, así que
%     se renombran una por una ---
\floatname{algorithm}{Algoritmo}
\renewcommand{\listalgorithmname}{Índice de algoritmos}

\algrenewcommand\algorithmicrequire{\textbf{Entrada:}}
\algrenewcommand\algorithmicensure{\textbf{Salida:}}
\algrenewcommand\algorithmicend{\textbf{fin}}
\algrenewcommand\algorithmicif{\textbf{si}}
\algrenewcommand\algorithmicthen{\textbf{entonces}}
\algrenewcommand\algorithmicelse{\textbf{si no}}
\algrenewcommand\algorithmicfor{\textbf{para}}
\algrenewcommand\algorithmicforall{\textbf{para todo}}
\algrenewcommand\algorithmicdo{\textbf{hacer}}
\algrenewcommand\algorithmicwhile{\textbf{mientras}}
\algrenewcommand\algorithmicrepeat{\textbf{repetir}}
\algrenewcommand\algorithmicuntil{\textbf{hasta que}}
\algrenewcommand\algorithmicloop{\textbf{ciclo}}
\algrenewcommand\algorithmicprocedure{\textbf{procedimiento}}
\algrenewcommand\algorithmicfunction{\textbf{función}}
\algrenewcommand\algorithmicreturn{\textbf{devolver}}
\algrenewcomment[1]{\hfill{\color{comentario}\itshape // #1}}

% Las dos que algorithmicx no trae. El \ del final es obligatorio: sin
% el TeX se come el espacio y sale "hastan-1" en vez de "hasta n-1"
\algnewcommand\Hasta{\textbf{hasta}\ }
\algnewcommand\Intercambiar{\textbf{intercambiar}\ }

% --- abreviaturas que se usan a cada rato ---
\newcommand{\N}{\mathbb{N}}
\newcommand{\Z}{\mathbb{Z}}
\newcommand{\Q}{\mathbb{Q}}
\newcommand{\R}{\mathbb{R}}
\newcommand{\C}{\mathbb{C}}
\newcommand{\abs}[1]{\left\lvert #1 \right\rvert}
\newcommand{\norm}[1]{\left\lVert #1 \right\rVert}
\newcommand{\set}[1]{\left\{ #1 \right\}}
\newcommand{\Ord}[1]{\mathcal{O}\!\left(#1\right)}
% Mayuscula la notacion grande, minuscula la chica
\newcommand{\Omg}[1]{\Omega\!\left(#1\right)}
\newcommand{\Tht}[1]{\Theta\!\left(#1\right)}
\newcommand{\ord}[1]{o\!\left(#1\right)}
\newcommand{\omg}[1]{\omega\!\left(#1\right)}
% Los tres limites sobre n, que salen en cada renglon del criterio
% del cociente para la notacion asintotica
\newcommand{\LM}{\lim_{n \to \infty}}
\newcommand{\LS}{\limsup_{n \to \infty}}
\newcommand{\LI}{\liminf_{n \to \infty}}

% --- encabezado de cada sesión: título a la izquierda, fecha a la
%     derecha. El argumento corto es el que sale en el índice ---
\newcommand{\clase}[2]{%
  \section[#1]{#1 \hfill \normalfont\small\textit{#2}}%
}

% --- portadilla y tabla de contenidos. Se llama una vez, justo
%     despues de \begin{document}. El renglon del profesor va dentro
%     del center, asi que \profesor puede traer varios nombres
%     separados por \\ y salen apilados y centrados sin tocar nada ---
\newcommand{\portadilla}{%
  \begin{center}
    {\LARGE\bfseries\materia\par}
    \vspace{0.4em}
    {\large \profesor \par}
    \vspace{0.2em}
    {\small \periodo\quad-\quad\alumno\par}
    \vspace{0.3em}
    \rule{\linewidth}{0.5pt}
  \end{center}
  \vspace{0.5em}
  \tableofcontents
  \vspace{1em}
  \rule{\linewidth}{0.5pt}%
}

% =====================================================================
%  DOCUMENTOS DE SOLUCIONES
%  Lo que usa una tanda de ejercicios resueltos. Una libreta no toca
%  nada de aqui, pero vive en el preambulo compartido para no tener
%  dos archivos que arreglar cuando cambie el estilo.
% =====================================================================

% --- encabezado compacto: titulo, subtitulo y regla. Sustituye a
%     \portadilla cuando el documento es una tanda de ejercicios: sin
%     indice, porque se hojea de corrido, no se navega ---
\newcommand{\encabezadoSoluciones}[2]{%
  % El paquete algorithm se carga con la opcion [section] para que en
  % una libreta el pseudocodigo se numere por clase. Aqui no hay
  % secciones, asi que esa numeracion sale como "Algoritmo 0.1". Se
  % corrige al vuelo, y solo en los documentos que llaman a este
  % encabezado: la libreta no se entera.
  \renewcommand{\thealgorithm}{\arabic{algorithm}}%
  \begin{center}
    {\LARGE\bfseries #1\par}
    \vspace{0.3em}
    {\small\itshape #2\par}
    \vspace{0.5em}
    \rule{\linewidth}{0.5pt}
  \end{center}
  \vspace{0.3em}%
}

% --- hoja de un solo ejercicio: un renglon chico arriba con el
%     titulo y su numero, y ya. Sin titulo grande y sin pie de pagina,
%     porque en una hoja de dos paginas cada centimetro que se va en
%     adornos es un paso del desarrollo que no cabe ---
\newcommand{\encabezadoEjercicio}[3]{%
  % Misma correccion que en \encabezadoSoluciones: sin secciones, el
  % pseudocodigo saldria numerado "0.1"
  \renewcommand{\thealgorithm}{\arabic{algorithm}}%
  \fancyhf{}%
  \fancyhead[L]{\small\textsc{#1} (ejercicio #2 de #3)}%
  % El nombre va a la derecha porque las hojas de respuestas se
  % entregan y el profesor pide que cada una lo lleve
  \fancyhead[R]{\small\alumno}%
  \renewcommand{\headrulewidth}{0.4pt}%
  \pagestyle{fancy}%
  \thispagestyle{fancy}%
}

% --- bloque tematico, numerado con romanos: I, II, III. Agrupa los
%     ejercicios que comparten tema. No dibuja regla: la del
%     \problema que sigue ya separa el titulo del primer ejercicio,
%     y dos rayas seguidas dejan un hueco feo ---
\newcounter{parte}
\renewcommand{\theparte}{\Roman{parte}}
\newcommand{\parte}[1]{%
  \bigskip
  \refstepcounter{parte}%
  \noindent{\large\bfseries \theparte. #1\par}%
}

% --- un ejercicio resuelto. Se numera solo y corrido a lo largo del
%     documento, sin reiniciar en cada parte, para poder decir "el 14"
%     sin aclarar de que bloque ---
\newcounter{problema}
\newcommand{\problema}[1]{%
  \bigskip
  \noindent\rule{\linewidth}{0.4pt}\par
  \vspace{0.3em}
  \refstepcounter{problema}%
  \noindent{\bfseries Ejercicio \theproblema\ --- #1\par}
  \vspace{0.3em}%
}

% --- rotulo del paso que sigue: Caso base, Hipotesis inductiva, Paso
%     inductivo. El texto sigue en el mismo renglon ---
\newcommand{\paso}[1]{\par\vspace{0.3em}\noindent\textbf{#1}\ }

% --- la respuesta de un ejercicio, centrada y en un cuadro. El
%     argumento va en modo matematico, sin los delimitadores ---
\newcommand{\respuesta}[1]{%
  \[
    \boxed{\;#1\;}
  \]
}

% --- resaltado inline, como pasarle el plumon a media frase. No
%     admite matematicas adentro: soul rompe con $...$, asi que una
%     formula que hay que destacar va en \respuesta ---
\sethlcolor{marcador}
\newcommand{\resaltar}[1]{\hl{#1}}

% --- justificacion al final de un renglon de \align, en chiquito y
%     entre parentesis, para decir de donde salio el paso ---
\newcommand{\razon}[1]{\quad\text{\small\itshape (#1)}}

%############################ END OF PREAMBULO.TEX #############################
%###############################################################################


\begin{document}

\encabezadoEjercicio{Panqueques de tamaño 1 o 2 con MergeSort}{13}{20}

\noindent\textbf{Convenciones.}
Todos los logaritmos son en base $2$ y $n$ es potencia de $2$.

\noindent Los panqueques $P[1 \ldots n]$ se etiquetan de arriba hacia
abajo y la operación $\textsc{Flip}(i, j)$, con $1 \le i \le j \le n$,
invierte el orden de los panqueques de la posición $i$ a la $j$ en
tiempo $\Ord{1}$ y con un costo de $j - i$ unidades de energía. La
energía del algoritmo es la suma de las energías de sus flips.
Suponga que todos los panqueques tienen tamaño $1$ o $2$. Diseñe un
algoritmo con tiempo de ejecución $\Ord{n \log n}$ que ordene los
panqueques utilizando $\Ord{n \log n}$ unidades de energía. Analice el
tiempo de ejecución y el consumo de energía de su algoritmo, y
demuestre su correctitud.

\paso{Algoritmo:} MergeSort donde dos mitades ordenadas
$1^{a} 2^{b} 1^{c} 2^{d}$ se mezclan con un solo flip sobre
$2^{b} 1^{c}$. Cada llamada devuelve cuántos unos tiene su segmento,
así que el flip se ubica sin recorrer nada.

\begin{algorithm}[H]
  \caption{Ordena una pila de panqueques de tamaño $1$ o $2$}
  \label{ej13:ej13-alg-ordena}
  \begin{algorithmic}[1]
    \Require $P[1 \ldots n]$ con $P[k] \in \{1, 2\}$, índices
             $1 \le p \le r \le n$
    \Ensure $P[p \ldots r]$ ordenado; devuelve su número de unos
    \Function{Ordena}{$P, p, r$}
      \If{$p = r$}
        \If{$P[p] = 1$}
          \State \Return $1$
        \Else
          \State \Return $0$
        \EndIf
      \EndIf
      \State $q \gets \lfloor (p + r) / 2 \rfloor$
      \State $a \gets \Call{Ordena}{P, p, q}$
        \Comment{unos de la mitad de arriba}
      \State $c \gets \Call{Ordena}{P, q + 1, r}$
        \Comment{unos de la mitad de abajo}
      \State $b \gets (q - p + 1) - a$
        \Comment{doses de la mitad de arriba}
      \If{$b > 0$ \textbf{y} $c > 0$}
        \State \Call{Flip}{$p + a,\; q + c$}
          \Comment{invierte $2^{b} 1^{c}$}
      \EndIf
      \State \Return $a + c$
    \EndFunction
  \end{algorithmic}
\end{algorithm}

\noindent La llamada inicial es $\textsc{Ordena}(P, 1, n)$.

\paso{Tiempo de ejecución:} fuera de las llamadas recursivas cada
invocación hace trabajo constante y a lo más un flip de costo
$\Ord{1}$:
\[
  T(n) = 2\,T(n/2) + \Tht{1}, \qquad T(1) = \Tht{1},
\]
con dos subproblemas de tamaño $n/2$ y $f(n) \in \Tht{1}$.
\begin{align*}
  n^{\log_2 2} &= n \razon{$\log_2 2 = 1$} \\
  f(n) \in \Tht{1} &\subseteq \Ord{n^{0}}
          \razon{el trabajo propio es constante} \\
          &= \Ord{n^{1 - \varepsilon}}
          \razon{con $\varepsilon = 1 > 0$}
\end{align*}
Aplica el caso 1, que concluye
$T(n) \in \Tht{n^{\log_2 2}} = \Tht{n}$.

\paso{Energía:} sobre un segmento de tamaño $m = r - p + 1$ el flip
de $i = p + a$ a $j = q + c$ cuesta
\begin{align*}
  j - i &= (q + c) - (p + a) \razon{posiciones del flip} \\
        &= \left[(q - p + 1) - a\right] + c - 1
           \razon{se suma y se resta $1$} \\
        &= b + c - 1 \razon{$b = (q - p + 1) - a$} \\
        &\le m - 1
           \razon{$b$ y $c$ cuentan panqueques distintos de
                  $P[p \ldots r]$}
\end{align*}
Sumando sobre las llamadas,
\[
  E(n) \le 2\,E(n/2) + n - 1, \qquad E(1) = 0,
\]
y se prueba por inducción que $E(n) \le n \log n$.

\paso{Caso base:} con $n = 1$,
\[
  E(1) = 0 = 1 \cdot \log 1.
\]

\paso{Hipótesis inductiva:} $E(n/2) \le \dfrac{n}{2} \log \dfrac{n}{2}$.

\paso{Paso inductivo:}
\begin{align*}
  E(n) &\le 2\,E(n/2) + n - 1 \razon{recurrencia de la energía} \\
       &\le 2 \cdot \frac{n}{2} \log \frac{n}{2} + n - 1
          \razon{hipótesis inductiva} \\
       &= n \log \frac{n}{2} + n - 1 \razon{se simplifica el $2$} \\
       &= n (\log n - 1) + n - 1
          \razon{$\log(n/2) = \log n - 1$} \\
       &= n \log n - n + n - 1 \razon{se distribuye $n$} \\
       &= n \log n - 1 \\
       &\le n \log n
\end{align*}

\paso{Correctitud:} por inducción sobre $m = r - p + 1$:
$\textsc{Ordena}(P, p, r)$ deja $P[p \ldots r]$ no decreciente, solo
permuta panqueques dentro de $[p, r]$ y devuelve su número de unos.

\paso{Caso base:} con $m = 1$ el segmento ya está ordenado, no hay
flip y se devuelve $1$ si $P[p] = 1$ y $0$ si $P[p] = 2$.

\paso{Hipótesis inductiva:} vale para todo segmento de tamaño menor
que $m$, con $m \ge 2$.

\paso{Paso inductivo:} como $p \le q < r$, ambas mitades son no vacías
y menores que $m$; por la hipótesis inductiva quedan
$P[p \ldots q] = 1^{a} 2^{b}$ y $P[q + 1 \ldots r] = 1^{c} 2^{d}$.
El único descenso de $1^{a} 2^{b} 1^{c} 2^{d}$ está entre $2^{b}$ y
$1^{c}$, y existe si y solo si $b > 0$ y $c > 0$. Esos bloques ocupan
las posiciones $p + a$ a $q + c$, así que el flip invierte
exactamente $2^{b} 1^{c}$:
\begin{align*}
  \left(2^{b} 1^{c}\right)^{R}
    &= \left(1^{c}\right)^{R} \left(2^{b}\right)^{R}
       \razon{invertir intercambia el orden de los bloques} \\
    &= 1^{c} 2^{b}
       \razon{un bloque constante invertido es él mismo}
\end{align*}
El segmento queda $1^{a+c} 2^{b+d}$, igual que sin flip cuando $b = 0$
o $c = 0$, y $a + c$ es su número de unos. El flip cae dentro de
$[p, r]$, así que no toca nada fuera del segmento.

\respuesta{T(n) \in \Tht{n} \subseteq \Ord{n \log n},
  \qquad E(n) \le n \log n \in \Ord{n \log n}}

\end{document}

%################### END OF EJERCICIO13PANQUEQUESUNODOS.TEX ####################
%###############################################################################
